\item Un líquido tiene una densidad $\rho$. (a) Demuestre que el cambio fraccionario en densidad para un cambio en temperatura $\Delta T$ es $\Delta \rho / \rho = -\beta \Delta T$. ¿Qué implica el signo negativo? (b) El agua pura tiene una densidad máxima de $1.0000\text{ g/cm}^3$ a 4.0 °C. A 10.0 °C, su densidad es $0.9997\text{ g/cm}^3$. ¿Cuál es $\beta$ para el agua en el rango de temperatura de 0 a 4.00 °C?
